\begin{exercisechunk}
\item \textbf{A small target sample can select the head.}
\label[exercise]{ex:l13-label-reversal}
\exerciseprofile{3}{2}{3}{\exproof\quad\excalculation\quad\exinterpretation}{%
Sharpen the generic target-ERM bound when only the orientation of a shared
Boolean feature remains to be learned.}
Use the label-reversal model of \cref{eq:label-reversal-model}.
Suppose the source procedure has an event of probability at least
$1-\delta_{\rm src}$ on which
$r:=A_0(\widehat h)\le\varepsilon\le1/4$.
Freeze $\widehat h$ and select one of the two heads by ERM on $m$
fresh target examples independent of all source data.
\begin{enumerate}[(a),beginpenalty=10000]
\item Conditional on a source sample in that event, prove that the
probability of choosing the higher-risk head is at most
$\exp(-m/8)$. Allow arbitrary tie breaking.
\item Deduce a guarantee that the final target risk is at most $\varepsilon$
with probability at least $1-\delta_{\rm src}-\delta_{\rm trg}$.
Give sufficient $n$ and $m$ when the source learner is joint ERM with
$n$ examples from each of $T$ tasks.
\item For fixed $\varepsilon\in(0,1/4)$ and fixed confidence, compare the
target sample requirement with the boxes lower bound for a learner using
only target data. Explain which information the source data supply and which
information the target data supply.
\end{enumerate}
\begin{hints}
For each target example, the two heads' mistake indicators sum to one.
Apply Hoeffding's inequality to the empirical risk of the better head.
\end{hints}
\end{exercisechunk}
